\documentclass[11pt]{amsart}
\usepackage[margin=1in]{geometry}
\usepackage{amsmath,amssymb,amsthm}
\usepackage{enumitem}
\usepackage[hidelinks]{hyperref}
\setlist{nosep,leftmargin=*}

\newtheorem{theorem}{Theorem}
\newtheorem{lemma}[theorem]{Lemma}
\newtheorem{proposition}[theorem]{Proposition}
\newtheorem{corollary}[theorem]{Corollary}
\theoremstyle{remark}
\newtheorem*{remark}{Remark}

\newcommand{\hX}{\widehat{X}}
\newcommand{\cross}{\textstyle\prod^{\times}}
\newcommand{\poch}[2]{(#1;t)_{#2}}

\title[The $e_k\star e_r$ Pieri rule for all $k$]{The $e_k\star e_r$ Pieri rule for all $k$: parabolic kernel, chain factorization, and a one-index $q$-series identity}
\author{Rick}
\date{2026-09-26}

\begin{document}

\begin{abstract}
For every $k\ge1$, $r\ge0$ and every number of variables $m\ge0$, we prove
\[
t^{-\binom k2}e_k(Y)\bullet e_r=\sum_{b=0}^{k}s^b\,F_{k-b}(t^{r-b})\,e_b\,e_{r+k-b},\qquad s=q^{-1},
\]
with the explicit polynomial $F_n$ defined in \eqref{eq:c} below. Via Hikita's Definition~3.4 and Lemma~3.3 (arXiv:2503.23597), the left side is $e_k\star e_r$. The cases $k=1$ and $k=2$ recover Hikita's Theorem~3.12 and our $W_r$. The cases $k=3$ and $k=4$ recover the Day~193 and Day~195 closed forms, which are now proved for all $r$.
\end{abstract}

\maketitle

\begin{center}
\begin{tabular}{ll}
\textbf{Author:} & Rick\\
\textbf{Date:} & 2026-09-26\\
\textbf{Recipient:} & Clio\\
\textbf{Source:} & Corresponds to work-in-progress commit \texttt{7682bfd}\\
& {\footnotesize(\texttt{proofs/2026-09-26-day207b-ek-star-er-general-k-PROVED.md}, \texttt{scripts/day207b/})}\\
\end{tabular}
\end{center}

\section{Conventions and statement}

\paragraph{Operators.} The conventions are those of the W$_r$ note.
\[
T_iF=t\,s_iF+(t-1)X_{i+1}\frac{s_iF-F}{X_i-X_{i+1}},\qquad
\pi F=X_1F(X_2,\dots,X_m,sX_1),\]
\[Y_i=t^{m-i}T_{i-1}\cdots T_1\,\pi\,T_{m-1}^{-1}\cdots T_i^{-1}.
\]

\paragraph{Symmetric functions.}
\begin{itemize}
\item $a_{ij}=(X_i-tX_j)/(X_i-X_j)$ and $\cross_A=\prod_{i\in A,\,j\notin A}a_{ij}$.
\item $E(z)=\prod_i(1+X_iz)=\sum_n e_nz^n$.
\item $Q(y)=\sum_n q_ny^n=E(-ty)/E(-y)$.
\item $[n]=[n]_t$.
\end{itemize}

\paragraph{Coefficients.} For $0\le j\le n$ put
\begin{equation}\label{eq:c}
\alpha_j=\frac{\prod_{i=1}^{j}(s-t^i)}{\poch tj},\quad \alpha_{-1}=0,\qquad
c(n,j)=\frac{\poch s{n-j}}{\poch t{n-j}}\bigl(\alpha_j-s\,t^{n-j}\alpha_{j-1}\bigr),\qquad
F_n(w)=\sum_{j=0}^n c(n,j)\,w^j .
\end{equation}
Set $c(n,j)=0$ outside the range $0\le j\le n$.

\begin{theorem}\label{thm:main}
For all $m\ge0$, $k\ge1$ and $r\ge0$,
\[
\sum_{r\ge0} z^r\,t^{-\binom k2}e_k(Y)\bullet e_r
=\sum_{j+b\le k}s^b t^{-kj}c(k-b,j)\;e_b\,z^{b-k}E(t^jz).
\]
Taking the coefficient of $z^r$ gives
\[
t^{-\binom k2}e_k(Y)\bullet e_r=\sum_{b=0}^k s^bF_{k-b}(t^{r-b})\,e_b e_{r+k-b}.
\]
\end{theorem}

\section{Hecke facts}

The $T_i$ satisfy the braid relations and $(T_i-t)(T_i+1)=0$. Hence:
\begin{itemize}
\item $T_w$ is well defined on reduced words;
\item $T_uT_v=T_{uv}$ whenever $\ell(uv)=\ell(u)+\ell(v)$.
\end{itemize}

For $J\subseteq S$, let $W^J$ be the set of minimal left coset representatives, so that $w=w^Jw_J$ uniquely and length-additively.

\paragraph{(Chain).} If $K\subseteq J$, then $W^J\times(W_J)^K\to W^K$, $(u,v)\mapsto uv$, is a length-additive bijection.
\begin{itemize}
\item If $s\in K$, then $\ell(uvs)=\ell(u)+\ell(vs)=\ell(u)+\ell(v)+1$.
\item Conversely, if $x\in W^K$ but $x_Js<x_J$ for some $s\in K$, then $\ell(xs)<\ell(x)$, which is a contradiction.
\end{itemize}

\paragraph{Spherical vectors.} If $T_kH=tH$ for all $s_k\in J$, then $T_vH=t^{\ell(v)}H$ for every $v\in W_J$.

\paragraph{Shifting through $\pi$.} $\pi T_k=T_{k+1}\pi$ for $1\le k\le m-2$.

\paragraph{Minimal representatives.} Let $J_n=\langle s_i:i\ne n\rangle$. For $D=\{d_1<\dots<d_n\}$, the word
\[
w(D)=(s_{d_1-1}\cdots s_1)(s_{d_2-1}\cdots s_2)\cdots(s_{d_n-1}\cdots s_n)
\]
sends $l\mapsto d_l$. Its length is $\sum_l(d_l-l)=\ell(w^D)$. Hence it is a reduced word for the minimal representative $w^D$. We set $\sigma^{(k)}=\sum_{|D|=k}T_{w(D)}$.

\section{The parabolic kernel $(A_k)$}

\begin{lemma}[Key Lemma]
Let $1\le c<d_1<\dots<d_n\le m$, and let $H$ satisfy $T_kH=tH$ for all $k\neq n$. Then
\[
(T_c\cdots T_{m-1})^{-1}\,T_{w(D)}H=t^{\,n-(m-c)}\,T_{w(D-1)}H.
\]
\end{lemma}

\begin{proof}
Let $x=s_c\cdots s_{m-1}$. It sends $m\mapsto c$ and $i\mapsto i+1$ for $c\le i<m$. Let $y=w(D-1)$; then $y(m)=m$.
\begin{itemize}
\item \emph{Length.} $x$ is increasing on $[1,m-1]$, and $xy(m)=c$ is inverted with exactly the $m-c$ positions $i$ with $y(i)\ge c$. So $\ell(xy)=\ell(y)+m-c$.
\item \emph{Coset.} $xy([n])=x(D-1)=D$. Therefore $xy=w^Dv$ with $v\in W_{J_n}$ and $\ell(v)=m-c-n$.
\end{itemize}
Hence $T_xT_yH=T_{w^D}T_vH=t^{m-c-n}T_{w^D}H$.
\end{proof}

\begin{proposition}[$A_k$]
Let $F$ be symmetric and $b_1<\dots<b_k$. Then $Y_{b_1}\cdots Y_{b_k}F=t^{\binom k2}T_{w(b)}\pi^kF$. Consequently $e_k(Y)F=t^{\binom k2}\sigma^{(k)}\pi^kF$.
\end{proposition}

\begin{proof}
Induct on $k$. Put $H=\pi^nF$; it is symmetric in the head and in the tail. By the Key Lemma,
\[
Y_c\,T_{w(D)}H=t^{m-c}T_{c-1}\cdots T_1\,\pi\,(T_c\cdots T_{m-1})^{-1}T_{w(D)}H=t^{n}\,T_{c-1}\cdots T_1\,\pi\,T_{w(D-1)}H .
\]
Shifting through $\pi$ turns this into $t^nT_{w(\{c\}\cup D)}\pi^{n+1}F$. Finally, $\binom n2+n=\binom{n+1}2$.
\end{proof}

\section{The $k$-set kernel $(K_k)$}

Let $\sigma_{[l,m]}=\sum_{a=l}^mT_{a-1}\cdots T_l$ and $K_l=\langle s_{l+1},\dots,s_{m-1}\rangle$.

\paragraph{(C1) Two factorizations.} Applying (Chain) along $K_k\subset K_{k-1}\subset\dots\subset K_0$, and also along $K_k\subset J_k$ (where $(W_{J_k})^{K_k}=S_k$), gives
\[
\sum_{u\in W^{K_k}}T_u=\sigma_{[1,m]}\cdots\sigma_{[k,m]}=\sigma^{(k)}\sum_{v\in S_k}T_v .
\]
Consequently, if $G$ is symmetric in the head and in the tail, then $\sigma_{[1,m]}\cdots\sigma_{[k,m]}G=[k]!\,\sigma^{(k)}G$.

\paragraph{(C2) Ordered formula.} Suppose $G$ is symmetric in the tail only. Iterating Day~205b Lemma~1, applied first in $X_2,\dots,X_m$ over $\mathbb Q(s,t)(X_1)$ and then at the top level, gives
\[
\sigma_{[1,m]}\cdots\sigma_{[k,m]}G=\sum_{(a_1,\dots,a_k)\text{ distinct}}G(X_{a_1},\dots,X_{a_k};\text{rest})\prod_{l=1}^k\prod_{j\notin\{a_1,\dots,a_l\}}a_{a_lj}.
\]

\paragraph{(C3) Poincar\'e identity.} Taking $m=k$ and $G=1$ in (C1) and (C2) gives $\sum_{w\in S_k}\prod_{l<l'}a_{w_lw_{l'}}=[k]!$.

\begin{proposition}[$K_k$]
If $G$ is symmetric in the head and in the tail, then $\sigma^{(k)}G=\sum_{|A|=k}G^{(A)}\cross_A$.
\end{proposition}

\begin{proof}
In (C2), group the ordered tuples by their underlying set $A$. The weight of each tuple factors as $\cross_A\cdot\prod_{l<l'}a_{a_la_{l'}}$. Summing over the orderings of $A$ gives $[k]!\,\cross_A$ by (C3). Now compare with (C1).
\end{proof}

\section{The generating function $(E_k)$}

Let $\varphi(u)=u(1+suz)/(1+uz)$. By $(A_k)$ and $(K_k)$,
\[
\Gamma_k:=\sum_rz^rt^{-\binom k2}e_k(Y)\bullet e_r=E(z)\sum_{|A|=k}\cross_A\prod_{a\in A}\varphi(X_a).
\]
This holds for all $m\ge0$; both sides vanish when $m<k$. Let $T_k$ denote the right side of Theorem~\ref{thm:main}.

\begin{lemma}[Recursion]
For $k,m\ge1$,
\[
[k]\,\Gamma_k(X)=\sum_iX_i(1+sX_iz)\prod_{j\ne i}a_{ij}\,\Gamma_{k-1}(\hX_i).
\]
\end{lemma}

\begin{proof}
Use (C2) and (C3), and peel off the first index $a_1=i$. Then use $\varphi(X_i)E(X;z)=X_i(1+sX_iz)E(\hX_i;z)$.
\end{proof}

\begin{lemma}[Step]\label{lem:step}
For every $m\ge0$ and $k\ge1$,
\[
\sum_iX_i(1+sX_iz)\prod_{j\ne i}a_{ij}\,T_{k-1}(\hX_i)=[k]\,T_k(X).
\]
\end{lemma}

\begin{proof}
Let $\gamma=t^jz$. Expand
\[
e_b(\hX_i)=\sum_c(-X_i)^ce_{b-c},\qquad E(\hX_i;\gamma)=\frac{E(\gamma)}{1+\gamma X_i},
\]
and split
\[
\frac{1+suz}{1+\gamma u}=st^{-j}+\frac{1-st^{-j}}{1+\gamma u}.
\]
Now apply Lemma~2 in the form $\sum_ig(X_i)\prod_{j\ne i}a_{ij}=(1-t)^{-1}\Omega[g]$, where $\Omega[u^n]=q_n$ for $n\ge1$. Two facts evaluate the result:
\begin{enumerate}[label=(\alph*)]
\item $\sum_c(-1)^ce_{b-c}q_{c+1}=(1-t^{b+1})e_{b+1}$;
\item $\displaystyle E(\gamma)\sum_c(-1)^ce_{b-c}\,\Omega\!\left[\frac{u^{1+c}}{1+\gamma u}\right]
=[w^b]\,\frac{E(w)E(t\gamma)-E(\gamma)E(tw)}{w-\gamma}$\\
$\displaystyle\qquad=-\sum_{n\le b}e_n\gamma^{n-1-b}\bigl(E(t\gamma)-t^nE(\gamma)\bigr).$
\end{enumerate}
For (b), first note that $\Omega[u/((1+wu)(1+\gamma u))]=(Q(-\gamma)-Q(-w))/(w-\gamma)$; then use $E(x)Q(-x)=E(tx)$.

Write $N'_{jb}=s^bt^{-(k-1)j}c(k-1-b,j)$ for the coefficients of $T_{k-1}$, and put $n=k-b'$. The coefficient of $e_{b'}z^{b'-k}E(t^jz)$ on the left side is then $(1-t)^{-1}s^{b'}t^{-kj}$ times
\[
(1-t^{b'})c(n,j)+t^{b'}\Bigl[(1-st^{-j})\sum_{p\ge0}(st^{-j})^pc(n-1-p,j)-t^n(1-st^{1-j})\sum_{p\ge0}(st^{1-j})^pc(n-1-p,j-1)\Bigr].
\]
By Lemma~\ref{lem:star} the bracket equals $(1-t^n)c(n,j)$. So the whole expression is $(1-t^k)c(n,j)$, which is the claim.
\end{proof}

\begin{lemma}[$\star$]\label{lem:star}
For $n\ge1$ and $j\ge0$,
\[
(1-t^n)c(n,j)=(1-st^{-j})\sum_{p\ge0}(st^{-j})^pc(n-1-p,j)-t^n(1-st^{1-j})\sum_{p\ge0}(st^{1-j})^pc(n-1-p,j-1).
\]
\end{lemma}

\begin{proof}
\emph{Generating functions.} Let $G(x)=\sum_n\frac{\poch sn}{\poch tn}x^n$. Comparing coefficients gives $(1-x)G(x)=(1-sx)G(tx)$. Let $C_j(x)=\sum_nc(n,j)x^n$. Then
\[
C_j(x)=x^j\bigl[\alpha_jG(x)-s\alpha_{j-1}G(tx)\bigr].
\]

\emph{Reformulation.} In terms of the $C_j$, the claim is
\[
\frac{C_j(x)(1-x)}{1-st^{-j}x}-C_j(tx)=-\frac{(1-st^{1-j})\,tx\,C_{j-1}(tx)}{1-st^{2-j}x}.
\]

\emph{Linear factorization.} Put $D_j=\alpha_j-s\alpha_{j-1}$. Since $\alpha_j=\alpha_{j-1}(s-t^j)/(1-t^j)$, we have
\[
\alpha_j(1-sy)-s\alpha_{j-1}(1-y)=D_j(1-st^{-j}y).
\]
Indeed, both sides are linear in $y$ and agree at $y=0$ and at $y=t^j/s$.

\emph{Reduction.} Substituting the factorization, both sides of the reformulated claim become the common factor $x^jG(t^2x)/(1-tx)$ times
\[
D_j(1-t^j)\qquad\text{and}\qquad t(s-t^{j-1})D_{j-1}
\]
respectively. These are equal, because $D_j=\alpha_{j-1}t^j(s-1)/(1-t^j)$ for $j\ge1$. For $j=0$ both sides vanish.
\end{proof}

\begin{proof}[Proof of Theorem~\ref{thm:main}]
Induct on $k$, for all $m$ simultaneously. The base case is $\Gamma_0=E(z)=T_0$. For $k\ge1$:
\begin{itemize}
\item if $m=0$, the Step Lemma gives $T_k=0=\Gamma_k$;
\item if $m\ge1$, combine the Recursion, the induction hypothesis in $m-1$ variables, and the Step Lemma.
\end{itemize}
Finally, $[z^r]\,z^{b-k}E(t^jz)=t^{j(r+k-b)}e_{r+k-b}$.
\end{proof}

\section{Remarks}

\paragraph{What is classical and what is new.} The kernel form of $e_k(Y)$ is classical in the $q$-shift DAHA; see Concha--Lapointe, arXiv:2307.02385, Lemmas~8 and~10. The level-one $\pi^k$ transfer and the $e$-basis Pieri rule for general $k$ are, as far as the Day~207 audit found, new.

\paragraph{Specializations.} The following are checked symbolically in $u=t^r$: $k=1$ gives Hikita's Theorem~3.12; $k=2$ gives $W_r$; $k=3,4$ give Days~193 and~195. The case $r=0$ recovers Hikita's Lemma~3.3, $e_k(Y)\bullet1=t^{\binom k2}e_k$; this is checked symbolically for $k\le7$.

\paragraph{Machine checks} (\texttt{scripts/day207b/}, all OK).
\begin{itemize}
\item Key Lemma and $(A_k)$: exact for $m\le6$.
\item (C2): checked.
\item Lemma~$\star$: checked for $n\le9$.
\item Step Lemma: checked with $z$ symbolic.
\item Full theorem against a direct AHA computation: 216 cases, including $m<k$.
\end{itemize}

\paragraph{Open.} The ${}_2\phi_1$ form $F_n(w)=(1-t^nw)\frac{\poch sn}{\poch tn}\,{}_2\phi_1(t^{1-n},t/s;t^{1-n}/s;t;w)$ is only computed ($n\le6$), and it is not needed for the theorem.

\end{document}
